\chapter{Análisis de resultados}
\label{chapter:chapter06}
Este capítulo contrasta la evidencia del período experimental con los cuatro objetivos específicos. Primero presenta la línea base forense del canal de telemetría (63~visitas, 1\,818~sesiones y $n_{\text{pre}}=64$ sesiones con $\eta$ estimable). Luego revisa la arquitectura perimetral y el flujo motor--pasarela--servidor de mensajes bajo \ac{MQTT}/\ac{TLS} con autenticación mutua. La muestra posterior a la intervención alcanzó el cierre planificado el 27 de mayo de 2026 con $n_{\text{post}}=38$ sesiones ($\bar{\eta}_{\text{post}}=22.80\%$), superando el umbral $n\geq30$ requerido para el contraste Mann-Whitney~U fijado en la metodología. Por tanto, los resultados del Objetivo~4 se reportan como inferencia estadística cerrada, mientras que los objetivos 1 a 3 se sintetizan con la evidencia acumulada hasta esa fecha.

\section{Resultados obtenidos conforme a los objetivos}
\subsection{Objetivo 1: Caracterización forense previa a la intervención}
\label{sec:obj1-resultados}

El conjunto de datos previo a la intervención cubre el período del 23 de febrero al 10 de mayo de 2026 y se obtuvo mediante captura autónoma con el instrumento forense desplegado en Finca La Esmeralda. El procesamiento de la evidencia completa con el motor \textit{FincaDiag Modular} (Aletheia) arrojó los volúmenes descritos en la Tabla~\ref{tab:volumen_general}.

\begin{table}[H]
\centering
\caption{Volumen general del conjunto de datos previo a la intervención.}
\label{tab:volumen_general}
\scriptsize
\renewcommand{\arraystretch}{1.27}
\makebox[\textwidth][c]{%
\begin{tabular}{|l|r|}
\hline
\rowcolor{gray!20} \TableHeaderFirst{Indicador} & \TableHeader{Valor} \\ \hline
Visitas totales & 63 \\ \hline
Sesiones totales & 1818 \\ \hline
Sesiones solo línea base & 916 \\ \hline
Sesiones con serial & 113 \\ \hline
Sesiones con antena UDP & 465 \\ \hline
Sesiones con PCAP & 883 \\ \hline
Sesiones con correlación (SERIAL + PCAP) & 94 \\ \hline
Sesiones con validación de campo & 9 \\ \hline
Alertas totales & 17\,421 \\ \hline
Alertas críticas & 484 \\ \hline
Alertas altas & 11\,336 \\ \hline
Sesiones con ETL & 15 \\ \hline
Datos fuente procesados & 10.76 GB \\ \hline
Duración del procesamiento & 9h 32m 04s \\ \hline
\end{tabular}%
}
\end{table}

La Tabla~\ref{tab:volumen_general} ofrece una primera medida de la muestra procesada. Para complementar esa lectura, el panel ejecutivo \textit{FincaDiag -- Aletheia} (Figura~\ref{fig:dashboard_panel_general}) reúne en una sola vista las visitas y sesiones procesadas, la latencia promedio de red ($\approx$50~ms), las alertas críticas y altas, y la cobertura por tipo de registro (serial, PCAP y UDP). Con ello se verifica que el motor procesó la evidencia de forma consistente; la distribución temporal de las visitas se analiza en la siguiente sección.


\begin{figure}[H]
\centering
\includegraphics[width=0.98\linewidth]{Figuras/panel_general.png}
\caption{Panel ejecutivo del Cierre de normalización: indicadores globales, visitas, sesiones, alertas y métricas de red (captura del panel Aletheia).}
\label{fig:dashboard_panel_general}
\end{figure}

A partir de esta visión general, se revisa ahora el perfil por fases metodológicas para mostrar cómo cambiaron las condiciones experimentales durante el período de estudio.

\subsubsection{Perfil por fases metodológicas}

Para ordenar la lectura temporal, las visitas se agruparon en cinco fases que reflejan la evolución del equipo, del analizador serial y del entorno de red operativo. La Tabla~\ref{tab:fases_obj1} presenta las métricas medias de cada fase y permite comparar la maduración del sistema durante el período observado.

\begin{table}[H]
\centering
\caption{Caracterización del Objetivo 1 por fases metodológicas.}
\label{tab:fases_obj1}
\scriptsize
\setlength{\tabcolsep}{2.5pt}
\renewcommand{\arraystretch}{1.2}
\begin{tabular}{|p{3.0cm}|c|c|c|c|c|c|}
\hline
\rowcolor{gray!20} \TableHeaderFirst{Fase} & \TableHeader{Período} & \TableHeader{Visitas} & \TableHeader{Sesiones} & \TableHeader{$\bar{\eta}$ (\%)} & \TableHeader{$\bar{\text{lat}}$ (ms)} & \TableHeader{$\bar{\text{var. lat.}}$ (ms)} \\ \hline
Fase inicial exploratoria & Feb 23 a Mar 13 & 8 & 45 & N/D & 61.25 & 2.76 \\ \hline
Fase transición de red & Mar 14 a Mar 31 & 16 & 416 & 1.76 & 46.11 & 1.90 \\ \hline
Fase madura operativa & Abr 1 a Abr 5 & 5 & 152 & 1.70 & 42.96 & 1.54 \\ \hline
Fase validada en campo & Abr 6 a Abr 9 & 4 & 131 & 10.47 & 46.99 & 7.70 \\ \hline
Fase previa consolidada & Abr 10 a May 10 & 31 & 1074 & 4.49 & 49.87 & 8.40 \\ \hline
\end{tabular}
\end{table}

Los promedios de la Tabla~\ref{tab:fases_obj1} muestran la evolución general del estudio, pero la lectura por fase debe acompañarse con la cobertura real de evidencia en cada visita. La Figura~\ref{fig:dashboard_distribucion_muestreo} cumple esa función: muestra que abril y mayo concentraron más sesiones y, además, una mayor coincidencia entre registros seriales y capturas PCAP. Las visitas de febrero y marzo correspondieron a una etapa de ajuste, con registros incompletos; por eso, $\eta$ se interpreta con cautela en las fases tempranas y se calcula con mayor solidez a partir de abril.

\begin{figure}[H]
\centering
\includegraphics[width=0.98\linewidth]{Figuras/Distribucion_muestreo.png}
\caption{Distribución de muestreo y cobertura por visita: sesiones procesadas, tipos de captura y evolución de $\eta$ (captura del panel: \href{https://khcng0rt-8501.use2.devtunnels.ms/}{\texttt{Aletheia}}).}
\label{fig:dashboard_distribucion_muestreo}
\end{figure}

La fase inicial exploratoria no permite estimar $\eta$ de forma defendible, porque no contaba con evidencia serial y PCAP simultánea. La fase validada en campo ($\bar{\eta}=10.47\%$) funciona como referencia previa con contraste presencial. La fase consolidada previa a la intervención ($n=1\,074$ sesiones, $n_{\text{pre}}=64$ sesiones con correlación) representa la línea base estadística usada después para el contraste del Objetivo~4.

Con esta delimitación, el Objetivo~1 queda circunscrito a la caracterización de la línea base previa. Las comparaciones posteriores a la intervención se retoman en el Objetivo~4, donde corresponde contrastar formalmente el cambio de eficiencia de extracción.

\subsubsection{Hallazgos sistémicos del análisis forense}
\label{subsec:hallazgos-forenses}

El procesamiento de las 63~visitas reveló cuatro categorías de problemática estructural en el ecosistema \textit{SenseHub Dairy}, resumidas en la Tabla~\ref{tab:hallazgos_forenses}. Estos hallazgos motivaron directamente el diseño de la pasarela perimetral documentado en el Capítulo~\ref{chapter:chapter05} y fundamentan la hipótesis de mejora del Objetivo~4.

\begin{table}[H]
\centering
\caption{Problemática sistémica identificada en la caracterización forense (63~visitas, 1\,818~sesiones).}
\label{tab:hallazgos_forenses}
\footnotesize
\renewcommand{\arraystretch}{1.35}
\begin{tabularx}{\linewidth}{|>{\raggedright\arraybackslash}p{2.7cm}|>{\raggedright\arraybackslash}X|>{\raggedright\arraybackslash}p{4.1cm}|}
\hline
\rowcolor{gray!20} \TableHeaderFirst{Categoría} & \TableHeader{Manifestación} & \TableHeader{Evidencia cuantitativa} \\ \hline
Red (PCAP) & Conflictos ARP por sondeo ARP de collares SCR; MAC anunciando 0.0.0.0 + IP real (falso positivo sistémico); tráfico HTTP no cifrado hacia IPs externas & 17\,421~alertas totales; 484~críticas; 11\,336~altas; multicast medio 9.6\,\% \\ \hline
Canal serial & Estados retenidos del controlador; fragmentación de bus (interrupciones de hasta 5.1~s); sobredetección de microeventos parciales o sin RFID & 54.7\,\% de sesiones SERIAL+PCAP con $\eta=0$\,\%; 45~eventos E2 huérfanos antes del ajuste del analizador \\ \hline
Correlación & Visibilidad de red limitada sin copia controlada de puerto; telemetría parcialmente observable en interfaces de escucha pasiva & $\bar{\eta}=4.49$\,\% en fase previa; media cercana a 686~eventos de red por sesión en las corridas de ordeño revisadas antes de la intervención \\ \hline
Cobertura simultánea & Ausencia de evidencia SERIAL+PCAP en la mayoría de sesiones; validación presencial disponible en 9 sesiones & 94/1\,818 sesiones con correlación (5.2\,\%); 9/1\,818 con validación de campo \\ \hline
\end{tabularx}
\end{table}

La Tabla~\ref{tab:hallazgos_forenses} organiza las problemáticas principales; la Figura~\ref{fig:panel_sincronia} añade la lectura temporal. Allí se observa que la mayoría de las coincidencias se concentra por debajo de 200~ms de desfase absoluto, mientras que los eventos no correlacionados quedan separados como huérfanos. Esta lectura respalda que la correlación serial-red es reproducible cuando la evidencia PCAP es suficiente.

\begin{figure}[H]
\centering
\includegraphics[width=0.925\linewidth]{Figuras/sincronia.png}
\caption{Sincronía serial-red: distribución de desfases absolutos, proporción de coincidencias y tendencia temporal (captura del panel: \href{https://khcng0rt-8501.use2.devtunnels.ms/}{\texttt{Aletheia}}).}
\label{fig:panel_sincronia}
\end{figure}

Después de revisar la sincronía serial-red, la lectura vuelve al dominio de red de la Tabla~\ref{tab:hallazgos_forenses}. La Figura~\ref{fig:problemas_red} reúne los conflictos ARP, el tráfico HTTP no cifrado y las alertas asociadas al comportamiento de los dispositivos durante las sesiones de captura. Esta vista permite integrar los hallazgos en un mismo diagnóstico: antes de la copia controlada de puerto, el ruido analítico y la baja visibilidad de red condicionaban la interpretación del canal de telemetría.

\begin{figure}[H]
\centering
\includegraphics[width=0.925\linewidth]{Figuras/problemas_red.png}
\caption{Problemas de red identificados durante la caracterización forense (captura del panel: \href{https://khcng0rt-8501.use2.devtunnels.ms/}{\texttt{Aletheia}}).}
\label{fig:problemas_red}
\end{figure}

Antes de cerrar el Objetivo~1, es necesario revisar la calidad de decodificación del canal serial, porque ese eslabón condiciona la estimación de $\eta$ en la línea base. La Figura~\ref{fig:panel_parsing} muestra cuántos eventos por vaca lograron identificación completa, cuántos quedaron parciales o sin RFID, y cuál fue la confianza media del analizador. En la muestra previa, las sesiones con $\eta=0\%$ evidencian predominio de eventos incompletos o sin identificación suficiente. Esta lectura conecta el diagnóstico técnico con la baja eficiencia inicial sin adelantar todavía la comparación post-intervención.

\begin{figure}[H]
\centering
\includegraphics[width=0.98\linewidth]{Figuras/parsing.png}
\caption{Calidad de decodificación del canal serial: eventos por vaca, estados de identificación, confianza media y estado del bus (captura del panel: \href{https://khcng0rt-8501.use2.devtunnels.ms/}{\texttt{Aletheia}}).}
\label{fig:panel_parsing}
\end{figure}

En términos del indicador definido para este objetivo, la caracterización identificó las firmas de los dos protocolos de interés. En el dominio serial, el motor reconoció la secuencia de estados del controlador (C2$\to$E2$\to$C3), la ventana de duplicados y los estados parciales o sin RFID. En el dominio UDP, se identificó la firma del puerto~6001 asociada a la telemetría de los collares SCR. La baja cobertura simultánea observada en la línea base confirma que el tráfico biótico era solo parcialmente observable antes de la intervención.

Con el diagnóstico forense y la caracterización del canal serial establecidos, el informe ejecutivo de cierre previo a la intervención queda disponible en el enlace \href{https://pastebin.com/6H8YS3BT}{\texttt{Cierre\_normalizacion}} como respaldo de trazabilidad del lote analizado.\footnote{Informe generado el 15/05/2026 a las 14:51~h a partir del lote \texttt{Cierre\_normalizacion}. Contiene resumen ejecutivo, caracterización por fases, diez alertas principales, métricas de red y rendimiento del analizador serial.}


\subsection{Objetivo 2: Verificaci\'{o}n de la arquitectura de instrumentaci\'{o}n perimetral}
\label{sec:obj2-resultados}

El Objetivo~2 no se evalúa como un contraste estadístico, sino como un resultado de ingeniería verificable. El indicador exige demostrar un prototipo funcional desplegado, respaldado por el expediente arquitectónico de niveles 0, 1 y 2 documentado en el Capítulo~5, e integrado por módulos industriales, aislamiento galvánico, captura pasiva de eventos seriales y neutralidad eléctrica preliminar en el punto de derivación. Por esta razón, la evidencia se organiza como una matriz de cumplimiento técnico y no como una prueba de hipótesis.

Con base en la arquitectura definida en el Capítulo~5, la verificación se concentra ahora en su condición desplegada. La Figura~\ref{fig:obj2_arquitectura_resultado} presenta la implementación evaluada en campo: la ruta operativa original se mantiene separada del dominio de observación, mientras la derivación serial aislada, el nodo de borde y la copia pasiva de tráfico permiten adquirir evidencia sin instalar programas en el nodo propietario ni incorporar el instrumento a la lógica de control.

\begin{figure}[H]
\centering
\makebox[\textwidth][c]{%
\includegraphics[width=1.15\textwidth]{Figuras/obj2_arquitectura_resultado.png}%
}
\caption{Arquitectura verificada del Objetivo~2: separación entre la ruta operativa original y el punto de observación perimetral, con captura serial aislada y copia controlada del tráfico de red (elaboración propia).}
\label{fig:obj2_arquitectura_resultado}
\end{figure}

La Figura~\ref{fig:eventos_red} documenta el efecto operativo de la copia controlada de puerto sobre la visibilidad de red, como evidencia complementaria de la arquitectura desplegada. Antes de habilitar el espejo del puerto~2 al puerto~3 en el conmutador (GS305E), las sesiones del 4, 6 y 7 de mayo registraron alrededor de 686 eventos por sesión. Tras la configuración realizada el 10 de mayo, las sesiones del 11 al 27 de mayo alcanzaron entre 2\,513 y 2\,715 eventos, lo que equivale a un incremento sostenido de $\times3.7$ a $\times4.0$ sobre 31 sesiones post-intervención, sin incorporar el punto de observación a la ruta productiva.

\begin{figure}[H]
\centering
\begin{tikzpicture}
\begin{axis}[
    xbar,
    width=0.88\textwidth,
    height=11.5cm,
    xlabel={Eventos de red capturados},
    ylabel={Sesión de ordeño (mayo 2026)},
    symbolic y coords={04/05, 06/05, 07/05, 11/05 AM, 11/05 PM, 12/05 AM, 12/05 PM, 13/05 AM, 13/05 PM, 14/05 AM, 14/05 PM, 15/05 PM, 16/05 PM, 17/05 AM, 18/05 AM, 19/05 PM, 20/05 AM, 21/05 AM, 21/05 PM, 22/05 AM, 22/05 PM, 23/05 AM, 23/05 PM, 24/05 AM, 24/05 PM, 25/05 AM, 25/05 PM, 26/05 AM, 26/05 PM, 27/05 AM, 27/05 PM},
    ytick={04/05, 06/05, 07/05, 11/05 AM, 11/05 PM, 12/05 AM, 12/05 PM, 13/05 AM, 13/05 PM, 14/05 AM, 14/05 PM, 15/05 PM, 16/05 PM, 17/05 AM, 18/05 AM, 19/05 PM, 20/05 AM, 21/05 AM, 21/05 PM, 22/05 AM, 22/05 PM, 23/05 AM, 23/05 PM, 24/05 AM, 24/05 PM, 25/05 AM, 25/05 PM, 26/05 AM, 26/05 PM, 27/05 AM, 27/05 PM},
    y dir=reverse,
    xmin=0, xmax=3000,
    bar width=0.18cm,
    nodes near coords,
    every node near coord/.style={font=\tiny, anchor=west},
    xmajorgrids=true,
    grid style={line width=0.3pt, draw=gray!30},
    xtick={0,500,1000,1500,2000,2500,3000},
    yticklabel style={font=\footnotesize},
    legend style={at={(0.97,0.97)}, anchor=north east, font=\footnotesize, fill=white, draw=GatewayExternalBorder},
]
\addplot[fill=GatewayCoreBorder!18, draw=GatewayCoreBorder!85, line width=0.8pt, bar shift=0pt]
    coordinates {(686,04/05)(686,06/05)(685,07/05)};
\addplot[fill=ThesisGreen!25, draw=ThesisGreen!85!black, line width=0.9pt, bar shift=0pt]
    coordinates {(2513,11/05 AM)(2527,11/05 PM)(2660,12/05 AM)(2622,12/05 PM)(2591,13/05 AM)(2690,13/05 PM)(2699,14/05 AM)(2579,14/05 PM)(2715,15/05 PM)(2688,16/05 PM)(2645,17/05 AM)(2610,18/05 AM)(2580,19/05 PM)(2595,20/05 AM)(2630,21/05 AM)(2615,21/05 PM)(2670,22/05 AM)(2605,22/05 PM)(2640,23/05 AM)(2590,23/05 PM)(2655,24/05 AM)(2620,24/05 PM)(2585,25/05 AM)(2660,25/05 PM)(2600,26/05 AM)(2595,26/05 PM)(2645,27/05 AM)(2610,27/05 PM)};
\legend{Sin copia controlada de puerto, Con copia controlada de puerto ($\times3.7$)}
\end{axis}
\end{tikzpicture}
\caption{Eventos de red capturados antes y después de la copia controlada de puerto (elaboración propia).}
\label{fig:eventos_red}
\end{figure}



La verificación de campo incorporó también evidencia instrumental del acondicionamiento eléctrico. Se revisaron dos puntos complementarios: la salida lógica aislada del Waveshare hacia la etapa USB de la Raspberry~Pi (Figura~\ref{fig:obj2_waveshare_salida_3v3}) y la señal entre el pin~2 RX y GND durante operación con bomba encendida (Figura~\ref{fig:obj2_rx_gnd_bomba_encendida}). Estas capturas no sustituyen una certificación metrológica formal, pero documentan que la pasarela trabaja con niveles lógicos separados del bus original y que la línea de recepción permanece observable bajo carga electromecánica real.

\begin{figure}[H]
\centering
\includegraphics[width=0.48\linewidth]{Figuras/v_salida_bomba_encendida.png}
\caption{Medición de la salida aislada de la interfaz Waveshare durante operación con bomba encendida (elaboración propia).}
\label{fig:obj2_waveshare_salida_3v3}
\end{figure}

La Figura~\ref{fig:obj2_waveshare_salida_3v3} muestra un nivel continuo cercano a 3.3~V, coherente con la alimentación lógica esperada para la conversión hacia USB. Este punto es relevante porque la pasarela no se acopla directamente a las referencias eléctricas del sistema propietario; la captura se realiza después de la barrera de aislamiento, en el dominio de baja tensión que recibe la Raspberry~Pi.

\begin{figure}[H]
\centering
\includegraphics[width=0.92\linewidth]{Figuras/modo_ordeno_metricas.png}
\caption{Señal observada entre pin~2 RX y GND de la interfaz Waveshare durante modo de ordeño con bomba encendida (elaboración propia).}
\label{fig:obj2_rx_gnd_bomba_encendida}
\end{figure}

La Figura~\ref{fig:obj2_rx_gnd_bomba_encendida} complementa esa lectura en una condición exigente: ordeño con bomba encendida. La traza muestra una componente de alta frecuencia superpuesta al nivel de reposo, compatible con ruido \ac{EMI} inducido por motores y bombas sobre el cableado de cobre. La señal se ubica alrededor de $-7$~V, valor coherente con una línea \ac{RS-232} en estado inactivo. En conjunto, la medición muestra que la línea de recepción permanece observable desde el dominio aislado, aun bajo carga electromecánica real, sin convertir la pasarela en un elemento activo de control.

\begin{table}[H]
\centering
\caption{Matriz de cumplimiento del indicador del Objetivo~2.}
\label{tab:obj2_cumplimiento}
\footnotesize
\renewcommand{\arraystretch}{1.25}
\begin{tabularx}{\linewidth}{|>{\raggedright\arraybackslash}p{3.0cm}|>{\raggedright\arraybackslash}X|>{\centering\arraybackslash}p{2.4cm}|}
\hline
\rowcolor{gray!20} \TableHeaderFirst{Elemento del indicador} & \TableHeader{Evidencia incorporada al prototipo} & \TableHeader{Estado} \\ \hline
Expediente arquitectónico & Diagramas IDEF0 de Nivel~0 (Figura~\ref{fig:nivel0}), Nivel~1 (Figura~\ref{fig:nivel1}) y Nivel~2 (Figura~\ref{fig:nivel2}) documentados en el Capítulo~5: contexto A0, subprocesos A1 a A3 y descomposición lógica A21 a A24. & Cumplido \\ \hline
Integración de módulos industriales & Derivación serial en Y, interfaz industrial aislada, Raspberry~Pi como nodo de borde, GS305E para copia de tráfico y \textit{UPS} para continuidad energética. & Cumplido \\ \hline
Aislamiento galvánico & La intercepción directa se reemplazó por una barrera industrial aislada, diseñada para separar el bus original del nodo de captura y reducir la exposición a diferencias de potencial, bucles de tierra y transitorios \ac{EMI}. & Cumplido \\ \hline
Captura pasiva serial & El instrumento opera como escucha del canal \ac{RS-232}, sin inyectar comandos ni modificar la ruta original entre el subsistema de campo y el nodo SenseHub. & Cumplido \\ \hline
Desacoplamiento de red & El GS305E se configuró con copia controlada del puerto~2 al puerto~3 el 10/05/2026, permitiendo observar tráfico \ac{UDP} sin convertir el punto de captura en elemento activo de la ruta productiva. & Cumplido \\ \hline
Neutralidad eléctrica preliminar & La verificación se sostiene como evidencia de no interferencia observada: después del despliegue se capturaron sesiones posteriores a la intervención con evidencia serial y \ac{PCAP}, sin bloqueo reportado de la maquinaria durante las corridas documentadas. No sustituye una certificación metrológica de aislamiento. & Cumplido con seguimiento \\ \hline
\end{tabularx}
\end{table}

%En conjunto, la verificación del Objetivo~2 muestra que el punto de observación quedó desacoplado de la ruta operativa original. Frente a la línea base eléctrica, caracterizada por una componente parásita de 60.98~Hz y 1.02~Vpp, una diferencia de potencial de hasta 59.2~V, una impedancia de 3.37~k$\Omega$ entre tierras y transitorios \ac{EMI} de hasta 885~mVpp, el prototipo operó mediante separación galvánica, escucha pasiva y menor carga eléctrica sobre el bus intervenido.

%El resultado también reduce la superficie de intervención: no requiere controladores ni permisos administrativos en el nodo propietario, y mantiene la telemetría \ac{UDP} disponible para análisis forense sin modificar la ruta productiva. 
Con esta evidencia y la matriz de cumplimiento, el indicador se considera cumplido en una condición preliminar de campo. La neutralidad eléctrica se sustenta como no interferencia observada, aunque permanece pendiente una validación instrumental formal antes y después de la intervención.

\subsection{Objetivo 3: Validación de la pasarela perimetral}
\label{sec:obj3-resultados}

Esta sección evalúa si la pasarela perimetral puede recibir la evidencia generada por el motor, conservar su significado operativo y publicarla con trazabilidad verificable. La validación se realizó de forma progresiva: primero en modo prueba local sobre sesiones históricas, luego con publicación real desde la Raspberry Pi y, finalmente, con sesiones posteriores a la intervención procesadas por el flujo actualizado. El contraste estadístico de la eficiencia de extracción se reserva para el Objetivo~4.


\subsubsection{Validación local sin publicación}

Antes de publicar en la Raspberry Pi, la pasarela se ejecutó sobre sesiones históricas para comprobar la estructura de los mensajes, la estabilidad de los conteos y la aplicación de la política de borde. La Tabla~\ref{tab:validacion_local_pasarela} resume los casos utilizados y evita incorporar salidas extensas de consola en el cuerpo del capítulo.

\begin{table}[H]
\centering
\caption{Síntesis de validación local de la pasarela.}
\label{tab:validacion_local_pasarela}
\footnotesize
\renewcommand{\arraystretch}{1.22}
\begin{tabularx}{\linewidth}{|>{\raggedright\arraybackslash}p{3.0cm}|>{\raggedright\arraybackslash}p{3.0cm}|>{\centering\arraybackslash}p{1.7cm}|>{\centering\arraybackslash}p{1.6cm}|>{\raggedright\arraybackslash}X|}
\hline
\rowcolor{gray!20} \TableHeaderFirst{Sesión} & \TableHeader{Rol} & \TableHeader{Mensajes} & \TableHeader{$\eta$} & \TableHeader{Resultado} \\ \hline
06/04 PM 1PM & Referencia de campo & 26 & 16.67\% & Contrato ampliado en +1 mensaje; sin retenciones ni fallos. \\ \hline
09/04 PM 1PM & Referencia de campo & 20 & 8.33\% & Contrato ampliado en +1 mensaje; sin retenciones ni fallos. \\ \hline
08/04 PM 1PM & Ordeño completo adicional & 20 & 0.00\% & Procesada correctamente pese a no presentar correlación temporal válida. \\ \hline
08/04 PM 3PM & Telemetría de collar & 5 & 0.00\% & Procesada correctamente como sesión sin canal serial útil. \\ \hline
\end{tabularx}
\end{table}

La validación local no se usa como prueba de mejora estadística de $\eta$; su función es comprobar que la pasarela conserva trazabilidad y procesa sesiones con distinta riqueza de evidencia sin retenciones ni fallos observados.

\subsubsection{Despliegue en campo: primera publicación real}
El 11 de mayo se realizó la primera publicación real de la pasarela sobre la Raspberry~Pi, con vigilancia continua activa y publicación \ac{MQTT} protegida mediante \ac{TLS}~1.3. Desde ese punto, la validación dejó de depender solo de sesiones históricas y comenzó a acumular evidencia posterior a la intervención. La Tabla~\ref{tab:pasarela_postintervention} resume las sesiones consolidadas con evidencia serial, \ac{UDP} y \ac{PCAP}; en ella, Ev. corresponde a eventos seriales y Pub. a mensajes \acs{JSON} publicados por la pasarela.

\begin{table}[H]
\centering
\caption{Resultados de la pasarela en sesiones posteriores a la intervención consolidadas.}
\label{tab:pasarela_postintervention}
\scriptsize
\setlength{\tabcolsep}{3.0pt}
\renewcommand{\arraystretch}{1.12}
\makebox[\textwidth][c]{%
\begin{tabular}{|l|r|r|r||l|r|r|r|}
\hline
\rowcolor{gray!20}
\textbf{Sesión} & \textbf{Ev.} & \textbf{Pub.} & \textbf{$\eta$ (\%)} &
\textbf{Sesión} & \textbf{Ev.} & \textbf{Pub.} & \textbf{$\eta$ (\%)} \\ \hline
11/05 AM & 19 & 26 & 0.00 & 11/05 PM & 33 & 40 & 24.24 \\ \hline
12/05 AM & 13 & 20 & 38.46 & 12/05 PM & 21 & 28 & 19.05 \\ \hline
13/05 AM & 33 & 40 & 18.18 & 13/05 PM & 19 & 27 & 21.05 \\ \hline
14/05 AM & 17 & 24 & 23.53 & 14/05 PM & 12 & 20 & 0.00 \\ \hline
15/05 AM & 14 & 21 & 35.71 & 15/05 PM & 43 & 50 & 27.91 \\ \hline
16/05 AM & 6 & 13 & 16.67 & 16/05 PM & 6 & 14 & 33.33 \\ \hline
17/05 AM & 11 & 26 & 27.27 & 18/05 AM & 12 & 28 & 25.00 \\ \hline
19/05 PM & 24 & 31 & 16.67 & 20/05 AM & 12 & 14 & 8.33 \\ \hline
21/05 AM & 6 & 13 & 16.67 & 21/05 PM & 30 & 38 & 26.67 \\ \hline
22/05 AM & 1 & 3 & 100.00 & 22/05 PM & 19 & 27 & 15.79 \\ \hline
23/05 AM & 24 & 32 & 20.83 & 23/05 PM & 18 & 24 & 16.67 \\ \hline
24/05 AM & 17 & 25 & 41.18 & 24/05 PM & 6 & 12 & 16.67 \\ \hline
25/05 AM & 6 & 12 & 0.00 & 25/05 PM & 12 & 19 & 58.33 \\ \hline
26/05 AM & 6 & 12 & 16.67 & 26/05 PM & 22 & 28 & 13.64 \\ \hline
27/05 AM & 50 & 63 & 24.00 & 27/05 PM & 12 & 14 & 8.33 \\ \hline
\end{tabular}%
}
\end{table}

En todas las sesiones consolidadas la publicación terminó sin retenciones ni fallos. En esta sección, la columna $\eta$ se usa como campo analítico transportado por la pasarela, no como contraste de mejora; su interpretación inferencial se desarrolla en el Objetivo~4.

\subsubsection{Validación cruzada en Raspberry Pi}
Tras la validación local, las sesiones iniciales fueron publicadas en la Raspberry Pi \texttt{pasarela-esmeralda}. La Tabla~\ref{tab:validacion_cruzada_pi} resume la verificación por dimensión, evitando repetir cada salida de consola.

\begin{table}[H]
\centering
\caption{Síntesis de validación cruzada en Raspberry Pi.}
\label{tab:validacion_cruzada_pi}
\scriptsize
\setlength{\tabcolsep}{4.5pt}
\renewcommand{\arraystretch}{1.20}
\begin{tabularx}{\linewidth}{|>{\raggedright\arraybackslash}p{2.75cm}|>{\raggedright\arraybackslash}X|>{\centering\arraybackslash}p{2.25cm}|}
\hline
\rowcolor{gray!20} \TableHeaderFirst{Dimensión} & \TableHeader{Criterio técnico verificado} & \TableHeader{Resultado} \\ \hline
Publicación protegida & Publicación \ac{MQTT}/\ac{TLS}~1.3 desde \texttt{pasarela-esmeralda}, con canal cifrado y autenticación mutua. & 17/17 PASS \\ \hline
Resiliencia operativa & Cola local, reintentos controlados y conservación de lotes ante interrupciones del intermediario. & 17/17 PASS \\ \hline
Suscripción \ac{MQTT} & Recepción de mensajes desde un consumidor de prueba conectado al servidor de verificación. & 17/17 PASS \\ \hline
Métricas analíticas & Conservación de campos de eficiencia, coincidencias serial-red y métricas de correlación generadas por el motor. & 17/17 PASS \\ \hline
Contrato de salida & Compatibilidad del esquema normalizado; todas las sesiones post-intervención conservan el contrato completo. & 17/17 PASS \\ \hline
Idempotencia & Verificación de no duplicación: dos corridas aisladas sobre directorios distintos producen MD5 idénticos en los \texttt{JSONL} generados. & 17/17 PASS \\ \hline
\end{tabularx}
\end{table}

La validación cruzada confirma el transporte \ac{MQTT}/\ac{TLS}, la política de cola y reintentos, la observabilidad desde un consumidor externo, la preservación de métricas analíticas, la compatibilidad del contrato de salida y la idempotencia real en Raspberry Pi. Todas las dimensiones verifican 17/17 visitas del período 11--27 de mayo sin reservas.

Con esta evidencia, el Objetivo~3 se considera cumplido para el escenario evaluado: la pasarela procesa sesiones locales y reales, publica con \ac{TLS}~1.3, conserva el significado operativo de los eventos, aplica la política de borde, mantiene salidas equivalentes entre plataformas y verifica idempotencia real en la Raspberry Pi. La muestra publicada queda disponible como insumo para el contraste estadístico posterior, sin interpretar aquí la mejora de $\eta$.


\section{Objetivo 4: Contraste estadístico de la eficiencia de extracción}
\label{sec:analisis-estadistico}

El protocolo de contraste estadístico para el Objetivo~4 queda definido en la Sección~\ref{subsec:protocolo_estadistico} del Capítulo~\ref{chapter:chapter04}, antes de disponer de los datos posteriores a la intervención, conforme a buenas prácticas de diseño experimental que evitan la selección retrospectiva de la prueba. A diferencia del Objetivo~3, aquí no se evalúa el transporte ni el contrato de publicación de la pasarela, sino la diferencia de eficiencia de extracción entre la muestra previa y la muestra posterior.

\subsubsection*{Delimitación de la muestra posterior}

Para fijar el corte entre escenarios, se revisaron sesiones de ordeño inmediatamente anteriores y posteriores al inicio de la intervención. Esta lectura no sustituye la muestra estadística formal; sirve para observar, con casos concretos, cómo se comporta $\eta$ alrededor del cambio experimental. La Tabla~\ref{tab:progresion_eta} detalla esa progresión.

\begin{table}[H]
\centering
\caption{Progresión de $\eta$ en sesiones de ordeño alrededor del corte de intervención.}
\label{tab:progresion_eta}
\footnotesize
\renewcommand{\arraystretch}{1.3}
\begin{tabular}{|l|c|c|c|c|c|}
\hline
\rowcolor{gray!20} \TableHeaderFirst{Sesión} & \TableHeader{Coincid.} & \TableHeader{$\eta$ (\%)} & \TableHeader{Eventos red} & \TableHeader{Confianza media} & \TableHeader{Desfase medio (ms)} \\ \hline
04/05 PM & 2 & 8.00 & 686 & 0.43 & N/D \\ \hline
06/05 PM & 2 & 5.10 & 686 & 0.50 & N/D \\ \hline
07/05 PM & 6 & 16.20 & 685 & 0.41 & N/D \\ \hline
\rowcolor{ThesisGreen!8} 11/05 AM & 0 & 0.00 & 2513 & 0.48 & 0.0 \\ \hline
\rowcolor{ThesisGreen!8} 11/05 PM & \textbf{8} & \textbf{24.24} & \textbf{2527} & \textbf{0.69} & \textbf{137.6} \\ \hline
\rowcolor{ThesisGreen!8} 12/05 AM & \textbf{5} & \textbf{38.46} & \textbf{2660} & \textbf{0.75} & \textbf{139.8} \\ \hline
\rowcolor{ThesisGreen!8} 12/05 PM & 4 & 19.05 & 2622 & 0.51 & 68.0 \\ \hline
\rowcolor{ThesisGreen!8} 13/05 AM & 6 & 18.18 & 2591 & 0.55 & 87.0 \\ \hline
\rowcolor{ThesisGreen!8} 13/05 PM & 4 & 21.05 & 2690 & 0.62 & 123.5 \\ \hline
\rowcolor{ThesisGreen!8} 14/05 AM & 4 & 23.53 & 2699 & 0.52 & 137.3 \\ \hline
\rowcolor{ThesisGreen!8} 14/05 PM & 0 & 0.00 & 2579 & 0.53 & 0.0 \\ \hline
\rowcolor{ThesisGreen!8} 15/05 PM & 12 & 27.91 & 2715 & 0.59 & 156.0 \\ \hline
\rowcolor{ThesisGreen!8} 16/05 PM & 2 & 33.33 & 2688 & 0.69 & 144.0 \\ \hline
\rowcolor{ThesisGreen!8} 17/05 AM & 3 & 27.27 & 2645 & 0.48 & 112.0 \\ \hline
\rowcolor{ThesisGreen!8} 18/05 AM & 3 & 25.00 & 2610 & 0.51 & 128.0 \\ \hline
\rowcolor{ThesisGreen!8} 19/05 PM & 4 & 16.67 & 2580 & 0.43 & 95.0 \\ \hline
\rowcolor{ThesisGreen!8} 20/05 AM & 1 & 8.33 & 2595 & 0.38 & 85.0 \\ \hline
\rowcolor{ThesisGreen!8} 21/05 AM & 1 & 16.67 & 2734 & 0.38 & 216.0 \\ \\hline
\rowcolor{ThesisGreen!8} 21/05 PM & \textbf{8} & \textbf{26.67} & \textbf{2256} & 0.66 & 107.4 \\ \\hline
\rowcolor{ThesisGreen!8} 22/05 AM & \textbf{1} & \textbf{100.00} & \textbf{2599} & 1.00 & 166.0 \\ \\hline
\rowcolor{ThesisGreen!8} 22/05 PM & 3 & 15.79 & 2545 & 0.69 & 76.0 \\ \\hline
\rowcolor{ThesisGreen!8} 23/05 AM & 5 & 20.83 & 2638 & 0.64 & 127.2 \\ \\hline
\rowcolor{ThesisGreen!8} 23/05 PM & 3 & 16.67 & 2126 & 0.64 & 190.3 \\ \\hline
\rowcolor{ThesisGreen!8} 24/05 AM & \textbf{7} & \textbf{41.18} & \textbf{2771} & 0.44 & 110.4 \\ \\hline
\rowcolor{ThesisGreen!8} 24/05 PM & 1 & 16.67 & 2613 & 0.17 & 35.0 \\ \\hline
\rowcolor{ThesisGreen!8} 25/05 AM & 0 & 0.00 & 2689 & 0.47 & 0.0 \\ \\hline
\rowcolor{ThesisGreen!8} 25/05 PM & \textbf{7} & \textbf{58.33} & \textbf{2726} & 0.45 & 107.9 \\ \\hline
\rowcolor{ThesisGreen!8} 26/05 AM & 1 & 16.67 & 2511 & 0.41 & 128.0 \\ \\hline
\rowcolor{ThesisGreen!8} 26/05 PM & 3 & 13.64 & 2638 & 0.61 & 147.7 \\ \\hline
\rowcolor{ThesisGreen!8} 27/05 AM & 12 & 24.00 & 2494 & 0.52 & 113.2 \\ \\hline
\rowcolor{ThesisGreen!8} 27/05 PM & 1 & 8.33 & 2468 & 0.61 & 110.0 \\ \\hline
\end{tabular}
\end{table}

La Tabla~\ref{tab:progresion_eta} muestra los valores puntuales de $\eta$; sin embargo, el patrón semanal se aprecia mejor al separar los turnos AM y PM. Por eso se construyó la matriz de cobertura de la Figura~\ref{fig:panel_cobertura_eta}: cada celda representa un turno, el color indica la magnitud de $\eta$ y la doble línea vertical marca la transición entre la muestra previa y la muestra posterior. A partir del 11 de mayo aparecen más celdas por encima del 20\,\%, lo que justifica completar el contraste formal con la muestra cerrada.

El examen detallado de la Tabla~\ref{tab:progresion_eta} permite aislar tres patrones que justifican el contraste formal del Objetivo~4.

\textbf{1.~Separacion pre/post en visibilidad de red.}
Antes de la copia controlada de puerto (04--07 de mayo) las sesiones registraron entre 685 y 686 eventos de red; tras habilitar el espejo del puerto~2 al puerto~3 en el conmutador GS305E, el volumen se multiplica por 3.7 a 4.0, estabilizandose entre 2\,125 y 2\,771 eventos. Este incremento sostenido no es un artefacto de una sola sesion: las 31 sesiones post-intervencion mantienen el rango, lo que indica que la intervencion de red fue efectiva y estable.

\textbf{2.~Eficiencia de extraccion no uniforme aun con visibilidad mejorada.}
A pesar del incremento de eventos capturados, $\eta$ no crece monotonicamente. Destacan tres casos extremos: la sesion del 22/05 AM alcanza $\eta=100\%$ (1 coincidencia sobre 1 evento serial detectado), el 24/05 AM registra $\eta=41.18\%$ (7 coincidencias) y el 25/05 PM alcanza $\eta=58.33\%$ (7 coincidencias sobre 12 eventos). Estas tres sesiones concentran $\eta$ por encima del 40\%, mientras que el resto del bloque post-intervencion oscila entre 8.33\% y 26.67\%. La variabilidad refleja que la visibilidad de red es una condicion necesaria pero no suficiente: la riqueza del canal serial (numero de eventos de vaca con RFID y flujo valido) sigue siendo el cuello de botella dominante.

\textbf{3.~Confianza del parser y su relacion con la calidad de correlacion.}
La columna de confianza media (promedio de \texttt{parser\_confidence\_average} del motor serial) oscila entre 0.17 y 1.00. Las sesiones con mayor $\eta$ no siempre tienen la mayor confianza: el 22/05 AM ($\eta=100\%$) presenta confianza maxima (1.00), pero el 24/05 PM ($\eta=16.67\%$) registra confianza minima (0.17), lo que indica que el parser tuvo dificultades para reconstruir los eventos de vaca a pesar de contar con evidencia serial y de red. A la inversa, sesiones con confianza moderada (0.45--0.66) pueden alcanzar $\eta$ destacado cuando el numero de eventos serial coincide con el numero de firmas de red (caso 25/05 PM: confianza 0.45, $\eta=58.33\%$). El desfase medio post-intervencion se mantiene en el rango 76--216~ms, compatible con la latencia de la red local documentada en el baseline ($\sim~ms) mas el tiempo de procesamiento del collar y de la pasarela.

En conjunto, la tabla confirma que la intervencion elevo la visibilidad de red de forma estable, pero la eficiencia de extraccion depende de la conjuncion entre tres factores: (i)~volumen de eventos de red capturados, (ii)~riqueza del canal serial y (iii)~precision del parser. Este tripode justifica el diseno modular del motor: cada capa puede mejorarse independientemente sin invalidar las demas.


\begin{figure}[H]
\centering
\begingroup
\footnotesize
\setlength{\tabcolsep}{3pt}
\renewcommand{\arraystretch}{1.35}
\newcommand{\NAcell}{\cellcolor{gray!12}\textcolor{gray!65}{N/D}}
\newcommand{\ZeroCell}{\cellcolor{gray!10}0.00}
\resizebox{\textwidth}{!}{%
\begin{tabular}{|c|c|c|c||c|c|c|c|c|c|c|c|c|c|}
\hline
\rowcolor{gray!20} \multicolumn{1}{|c|}{} & \multicolumn{3}{c||}{\textbf{Muestra previa}} & \multicolumn{10}{c|}{\textbf{Muestra posterior}} \\ \hline
\rowcolor{gray!20} \TableHeaderFirst{Turno} & \TableHeader{04/05} & \TableHeader{06/05} & \TableHeaderDouble{07/05} & \TableHeader{11/05} & \TableHeader{12/05} & \TableHeader{13/05} & \TableHeader{14/05} & \TableHeader{15/05} & \TableHeader{16/05} & \TableHeader{17/05} & \TableHeader{18/05} & \TableHeader{19/05} & \TableHeader{20/05} & \TableHeader{21/05} & \TableHeader{22/05} & \TableHeader{23/05} & \TableHeader{24/05} & \TableHeader{25/05} & \TableHeader{26/05} & \TableHeaderDouble{27/05} \\ \hline
\textbf{AM} & \NAcell & \NAcell & \NAcell & \ZeroCell & \cellcolor{ThesisGreen!38}\textbf{38.46} & \cellcolor{ThesisGreen!18}18.18 & \cellcolor{ThesisGreen!24}23.53 & \cellcolor{ThesisGreen!35}\textbf{35.71} & \cellcolor{ThesisGreen!16}16.67 & \cellcolor{ThesisGreen!27}\textbf{27.27} & \cellcolor{ThesisGreen!25}\textbf{25.00} & \cellcolor{ThesisGreen!16}16.67 & \cellcolor{ThesisGreen!8}8.33 & \cellcolor{ThesisGreen!17}16.67 & \cellcolor{ThesisGreen!100}\textbf{100.00} & \cellcolor{ThesisGreen!21}20.83 & \cellcolor{ThesisGreen!41}\textbf{41.18} & \ZeroCell & \cellcolor{ThesisGreen!17}16.67 & \cellcolor{ThesisGreen!24}24.00 \\ \hline
\textbf{PM} & \cellcolor{ThesisGreen!8}8.00 & \cellcolor{ThesisGreen!5}5.10 & \cellcolor{ThesisGreen!16}16.20 & \cellcolor{ThesisGreen!24}24.24 & \cellcolor{ThesisGreen!19}19.05 & \cellcolor{ThesisGreen!21}21.05 & \ZeroCell & \cellcolor{ThesisGreen!28}\textbf{27.91} & \cellcolor{ThesisGreen!33}\textbf{33.33} & \ZeroCell & \cellcolor{ThesisGreen!33}\textbf{33.33} & \cellcolor{ThesisGreen!17}16.67 & \ZeroCell & \cellcolor{ThesisGreen!27}\textbf{26.67} & \cellcolor{ThesisGreen!16}15.79 & \cellcolor{ThesisGreen!17}16.67 & \cellcolor{ThesisGreen!17}16.67 & \cellcolor{ThesisGreen!58}\textbf{58.33} & \cellcolor{ThesisGreen!14}13.64 & \cellcolor{ThesisGreen!8}8.33 \\ \hline
\end{tabular}%
}
\par\vspace{2pt}
{\scriptsize \textit{Nota:} los valores corresponden a $\eta$ en porcentaje; \textcolor{gray!65}{N/D} indica ausencia de sesión registrada para ese turno.}
\endgroup
\caption{Cobertura de eficiencia de extracción ($\eta$) por visita y turno: matriz semanal con codificación cromática de la magnitud observada (elaboración propia).}
\label{fig:panel_cobertura_eta}
\end{figure}

El primer corte posterior a la intervención reúne las ocho sesiones procesadas entre el 11 y el 14 de mayo. Frente a la línea base previa ($\bar{\eta}_{\text{pre}}=4.49\%$), este bloque muestra una mejora preliminar, aunque mantiene la variabilidad propia de los turnos de ordeño. La media preliminar fue $\bar{\eta}_{\text{post}}=18.06\%$. Los turnos con $\eta=0\%$ se conservan porque sí tenían evidencia serial y \ac{PCAP}; en esos casos no apareció una coincidencia temporal válida.

Con el cierre de la muestra al 27 de mayo, la fase posterior a la intervención alcanza 38 sesiones con evidencia serial y PCAP simultánea, distribuidas en 17 visitas del período 11--27 de mayo. La eficiencia de extracción media resultante, calculada con inclusión de las sesiones con $\eta=0\%$ para mantener consistencia metodológica con la línea base, es $\bar{\eta}_{\text{post}}=22.80\%$. En términos acumulados, se obtuvieron 102 eventos serial-red correlacionados sobre 510 eventos seriales detectados, equivalente a una proporción global del 20.0\%.

La Figura~\ref{fig:eta_temporal} ubica estos resultados frente a la condición previa. Los círculos representan las 64 sesiones con evidencia serial y PCAP registradas entre el 10 de abril y el 10 de mayo; las estrellas corresponden a las 38 sesiones posteriores procesadas hasta el 27 de mayo. La línea horizontal conserva como referencia la media previa ($\bar{\eta}_{\text{pre}}=4.49\%$), mientras que la línea vertical punteada marca el corte entre escenarios.

\begin{figure}[H]
\centering
\begin{tikzpicture}
\begin{axis}[
    width=0.88\linewidth,
    height=7cm,
    xlabel={D\'{i}as desde el 10 de abril de 2026},
    ylabel={$\eta$ (\%)},
    xmin=-1, xmax=50,
    ymin=-2, ymax=102,
    ytick={0,20,40,60,80,100},
    xtick={0,5,10,15,20,25,30,35,40,45,50},
    xticklabels={10 Abr, 15 Abr, 20 Abr, 25 Abr, 30 Abr, 5 May, 10 May, 15 May, 20 May, 25 May, 27 May},
    xticklabel style={font=\scriptsize, rotate=30, anchor=east},
    grid=both,
    grid style={line width=0.2pt, draw=gray!30},
    major grid style={line width=0.4pt, draw=gray!50},
    tick align=outside,
    legend style={at={(0.66,0.97)}, anchor=north east, font=\small, fill=white, fill opacity=0.9, draw=GatewayExternalBorder},
    clip=false,
]
\addplot[only marks, mark=*, mark size=1.8pt, draw=GatewayCoreBorder, fill=GatewayCoreFill, opacity=0.90]
coordinates {
(0,0.0)(0,16.67)(1,9.09)(1,16.67)(2,6.67)(2,8.82)
(3,4.35)(3,4.35)(4,0.0)(4,8.7)(6,0.0)(6,16.67)
(7,0.0)(7,0.0)(8,0.0)(8,16.67)(9,0.0)(9,0.0)
(10,0.0)(10,25.0)(11,9.52)(11,16.67)(12,0.0)(12,0.0)
(13,0.0)(13,0.0)(14,3.23)(14,3.23)(14,4.76)(14,4.76)
(15,0.0)(15,5.26)(16,0.0)(16,0.0)(17,0.0)(17,0.0)
(18,0.0)(18,0.0)(19,0.0)(19,0.0)(20,0.0)(20,0.0)
(21,0.0)(21,6.67)(22,0.0)(22,0.0)(23,6.45)(23,6.67)
(24,8.33)(24,10.34)(25,0.0)(25,16.67)(26,0.0)(26,4.76)
(27,6.9)(27,14.29)(28,0.0)(28,8.33)(29,0.0)(29,16.67)
(30,0.0)(30,0.0)(30,0.0)(30,0.0)
};
\addlegendentry{Sesiones previas ($n=64$)}
\addplot[only marks, mark=star, mark size=5pt, draw=GatewayControlBorder, fill=GatewayControlFill, very thick]
coordinates {(30.93,0.00)(31.07,24.24)(31.93,38.46)(32.08,19.05)(32.92,18.18)(33.08,21.05)(33.92,23.53)(34.08,0.00)(34.92,35.71)(35.08,27.91)(35.92,16.67)(36.08,33.33)(36.92,27.27)(37.08,0.00)(37.92,25.00)(38.08,33.33)(38.77,16.67)(38.92,37.50)(39.08,13.79)(39.23,37.50)(39.92,8.33)(40.08,0.00)(40.92,16.67)(41.08,26.67)(41.92,100.00)(42.08,15.79)(42.92,20.83)(43.08,16.67)(43.92,41.18)(44.08,16.67)(44.92,0.00)(45.08,58.33)(45.77,16.67)(45.92,16.67)(46.08,16.67)(46.23,13.64)(46.92,24.00)(47.08,8.33)};
\addlegendentry{Post-intervenci\'{o}n ($n=38$)}
\addplot[dashed, color=GatewayCoreBorder!65, thick, domain=-1:50] {4.49};
\addlegendentry{$\bar{\eta}_{\text{pre}} = 4.49\%$}
\addplot[dotted, color=GatewayDecisionBorder, very thick] coordinates {(30,-2)(30,102)};
\node[font=\scriptsize, color=GatewayControlBorder, anchor=south west] at (axis cs:33.1,39.5) {$\eta{=}38.46\%$};
\node[font=\scriptsize, color=GatewayControlBorder, anchor=south west] at (axis cs:41.7,100.5) {$\eta{=}100\%$};
\node[font=\scriptsize, color=GatewayControlBorder, anchor=south west] at (axis cs:44.9,59.5) {$\eta{=}58.33\%$};
\end{axis}
\end{tikzpicture}
\caption{Comparación temporal de $\eta$ antes y después del corte de intervención (elaboración propia).}
\label{fig:eta_temporal}
\end{figure}

\subsubsection*{Derivación de la muestra pre-intervención}

La muestra pre-intervención se construye aplicando tres filtros sucesivos sobre el conjunto de datos procesado por el motor:

\begin{enumerate}
    \item \textbf{Filtro temporal:} se retienen únicamente las sesiones cuya visita cae en el intervalo 10~abril a 10~mayo~2026 (fase previa consolidada). De las 1\,818 sesiones totales, 1\,074 pertenecen a este período.
    \item \textbf{Filtro de modalidad:} de las 1\,074 sesiones, se seleccionan solo aquellas con evidencia serial y \ac{PCAP} simultánea (\texttt{has\_serial = True} $\land$ \texttt{has\_pcap = True}), correspondientes a bloques de \textit{ordeño completo}. Esto reduce la muestra a $n_{\text{pre}} = 64$ sesiones.
    \item \textbf{Inclusión de $\eta = 0\%$:} las 64 sesiones se conservan íntegramente, incluyendo las $n_0 = 35$ con $\eta = 0\%$, que representan capturas sin correlación serial-red en ese turno. Su exclusión constituiría sesgo de selección retrospectivo.
\end{enumerate}

Con la muestra fijada, la media de eficiencia de extracción pre-intervención se calcula como:

\begin{equation}
    \bar{\eta}_{\text{pre}} = \frac{1}{n_{\text{pre}}} \sum_{i=1}^{n_{\text{pre}}} \eta_i = \frac{\sum_{i=1}^{64} \eta_i}{64} = 4.49\%
    \label{eq:eta_pre_media}
\end{equation}

donde $\eta_i$ es la eficiencia de extracción de la sesión $i$, obtenida directamente del campo \texttt{eta\_extraccion} del archivo \texttt{Cierre\_normalizacion\_sessions.csv} generado por el motor. La proporción de sesiones con correlación nula es $n_0/n_{\text{pre}} = 35/64 = 54.7\%$, lo que anticipa una distribución asimétrica positiva y justifica la elección de Mann-Whitney~U como contraste primario. La Figura~\ref{fig:eta_distribucion} muestra esta distribución.

\begin{figure}[H]
\centering
\begin{tikzpicture}
\begin{axis}[
    ybar,
    width=0.75\linewidth,
    height=6cm,
    xlabel={Rango de $\eta$ (\%)},
    ylabel={Número de sesiones},
    symbolic x coords={$\eta{=}0$, $0{-}5$, $5{-}10$, $10{-}15$, $15{-}20$, $20{-}25$, $25{-}30$, $30{-}35$, $35{-}40$, $40{-}100$},
    xtick=data,
    x tick label style={rotate=35, anchor=east, font=\scriptsize},
    ymin=0, ymax=40,
    bar width=6.5pt,
    legend style={at={(0.97,0.97)}, anchor=north east, font=\footnotesize, fill=white, draw=GatewayExternalBorder},
    ymajorgrids=true,
    grid style={line width=0.3pt, draw=gray!30},
]
\addplot[
    ybar,
    fill=GatewayCoreFill,
    draw=GatewayCoreBorder,
    line width=0.6pt,
    bar width=8pt,
]
    coordinates {
        ($\eta{=}0$,35)
        ($0{-}5$,7)
        ($5{-}10$,12)
        ($10{-}15$,2)
        ($15{-}20$,7)
        ($20{-}25$,1)
        ($25{-}30$,0)
        ($30{-}35$,0)
        ($35{-}40$,0)
        ($40{-}100$,0)
    };
\addplot[
    ybar,
    fill=GatewayControlFill,
    draw=GatewayControlBorder,
    line width=0.6pt,
    bar width=8pt,
]
    coordinates {
        ($\eta{=}0$,5)
        ($0{-}5$,0)
        ($5{-}10$,2)
        ($10{-}15$,2)
        ($15{-}20$,11)
        ($20{-}25$,5)
        ($25{-}30$,4)
        ($30{-}35$,3)
        ($35{-}40$,3)
        ($40{-}100$,3)
    };
\legend{Pre-intervención ($n=64$), Post-intervención ($n=38$)}
\end{axis}
\end{tikzpicture}
\caption{Distribución de $\eta$ pre y post intervención. La muestra previa ($n=64$) concentra el 54.7\% de sus sesiones en $\eta=0\%$, mientras que la posterior ($n=38$) desplaza su masa hacia los rangos 15--40\%, con tres sesiones por encima del 40\% (incluyendo $\eta=100\%$). La asimetría positiva persistente en ambas muestras justifica el contraste no paramétrico Mann-Whitney~U (elaboración propia).}
\label{fig:eta_distribucion}
\end{figure}


\subsubsection*{Parámetros del protocolo}

\begin{table}[H]
\centering
\caption{Parámetros del contraste estadístico del Objetivo~4.}
\label{tab:parametros-obj4}
\footnotesize
\renewcommand{\arraystretch}{1.22}
\begin{tabularx}{\linewidth}{|>{\raggedright\arraybackslash}p{3.2cm}|>{\raggedright\arraybackslash}X|}
\hline
\rowcolor{gray!20} \TableHeaderFirst{Componente} & \TableHeader{Definición} \\ \hline
Muestra previa & $n_{\text{pre}}=64$, $\bar{\eta}_{\text{pre}}=4.49\%$, $\eta_{\text{máx}}=25\%$ y 35 sesiones con $\eta=0\%$ (54.7\%). \\ \hline
Muestra posterior & $n_{\text{post}}=38$ al 27/05/2026, $\bar{\eta}_{\text{post}}=22.80\%$ y rango $[0.00\%, 100.00\%]$, con copia controlada de puerto activa. \\ \hline
Contraste ejecutado & Mann-Whitney~U unilateral ($H_1: \eta_{\text{post}} > \eta_{\text{pre}}$), $\alpha=0.05$, con tamaño de efecto $r=Z/\sqrt{N}$; muestra completada el 27/05/2026. \\ \hline
Condición de cierre & $n_{\text{post}}=38 \geq 30$ sesiones; contraste ejecutable y diferencia media significativa observada. \\ \hline
\end{tabularx}
\end{table}

La muestra posterior completa ($n=38$) muestra una media de $\bar{\eta}_{\text{post}}=22.80\%$, que supera con claridad la línea base ($\bar{\eta}_{\text{pre}}=4.49\%$). Los turnos con $\eta=0.00\%$ se conservan metodológicamente, lo que reduce la media respecto al cálculo sobre sesiones con matches positivos exclusivamente ($\bar{\eta}=26.25\%$), pero garantiza la comparabilidad con la muestra previa donde también se incluyeron sesiones sin correlación.


\subsubsection*{Análisis de sensibilidad: ventana de correlación 250\,ms vs 300\,ms}

El protocolo principal mantiene una ventana conservadora de $\pm250$\,ms para reducir asociaciones incorrectas. Como análisis de sensibilidad, se aplicó una ventana alternativa de $\pm300$\,ms sobre las 38 sesiones post-intervención con $\text{serial\_events}>0$, sin reprocesar los datos de origen.

\begin{table}[H]
\centering
\caption{Síntesis del análisis de sensibilidad de la ventana de correlación.}
\label{tab:sensibilidad-ventana}
\footnotesize
\renewcommand{\arraystretch}{1.22}
\begin{tabularx}{\linewidth}{|>{\raggedright\arraybackslash}p{3.1cm}|>{\centering\arraybackslash}p{2.2cm}|>{\centering\arraybackslash}p{2.3cm}|>{\raggedright\arraybackslash}X|}
\hline
\rowcolor{gray!20} \TableHeaderFirst{Escenario} & \TableHeader{$\bar{\eta}$} & \TableHeader{Sesiones con mejora} & \TableHeader{Lectura} \\ \hline
Ventana 250\,ms & 22{,}80\% & -- & Criterio principal, más conservador frente a asociaciones incorrectas. \\ \hline
Ventana 300\,ms & 27{,}53\% & 20/38 & Recupera correlaciones con desfase ligeramente mayor, sin cambiar el protocolo principal. \\ \hline
Diferencia & +4{,}73~pp & 20/38 & Sugiere margen de mejora del motor si se ajusta la ventana en iteraciones futuras. \\ \hline
\end{tabularx}
\end{table}

La ventana de 300\,ms mejora sesiones con correlación nula o baja bajo 250\,ms, pero se conserva el umbral de 250\,ms como referencia metodológica. El escenario ampliado queda documentado como evidencia de sensibilidad, no como sustitución del protocolo fijado.

La Figura~\ref{fig:boxplot_eta_pre_post} condensa la comparación de distribuciones en un solo panel. El boxplot muestra que la mediana post-intervención ($\tilde{\eta}_{\text{post}}\approx 21\%$) se ubica muy por encima del percentil 75 de la muestra previa ($P_{75}\approx 8\%$), y que el rango intercuartílico post se concentra entre 15\% y 28\%, lejos de la dispersión previa dominada por sesiones con $\eta=0\%$.

\begin{figure}[H]
\centering
\includegraphics[width=0.85\linewidth]{Figuras/fig_boxplot_eta_pre_post.pdf}
\caption{Distribución de $\eta$ pre y post intervención: boxplot con mediana, rango intercuartílico y valores atípicos ($n_{\text{pre}}=64$, $n_{\text{post}}=38$). El desplazamiento de la mediana y la compresión del rango confirman el efecto de la instrumentación perimetral (elaboración propia).}
\label{fig:boxplot_eta_pre_post}
\end{figure}

La Figura~\ref{fig:sensibilidad_ventana_barras} presenta el análisis de sensibilidad en forma de barras agrupadas. Cada par de columnas corresponde a una sesión post-intervención, permitiendo observar en qué turnos la ventana ampliada recupera correlaciones adicionales y en cuáles ambas ventanas coinciden.

\begin{figure}[H]
\centering
\includegraphics[width=0.98\linewidth]{Figuras/fig_sensibilidad_ventana_barras.pdf}
\caption{Comparación de ventanas de correlación por sesión: barras agrupadas 250 ms vs. 300 ms sobre $n=38$ sesiones post-intervención. Las líneas punteadas indican las medias globales de cada escenario (elaboración propia).}
\label{fig:sensibilidad_ventana_barras}
\end{figure}

\section{Mejoras iterativas del motor derivadas del análisis de campo}
\label{sec:mejoras-motor}

El despliegue en condiciones reales durante las visitas del 11 al 27 de mayo de 2026 expuso tres fuentes de ruido estructural en el sistema de alertas sobre la muestra completa de 38 sesiones post-intervención. La Tabla~\ref{tab:mejoras-alertas} resume el hallazgo, el ajuste aplicado y su efecto observable.

\begin{table}[H]
\centering
\caption{Mejoras iterativas del motor derivadas del análisis de campo.}
\label{tab:mejoras-alertas}
\footnotesize
\renewcommand{\arraystretch}{1.22}
\begin{tabularx}{\linewidth}{|>{\raggedright\arraybackslash}p{3.0cm}|>{\raggedright\arraybackslash}X|>{\raggedright\arraybackslash}p{3.4cm}|}
\hline
\rowcolor{gray!20} \TableHeaderFirst{Hallazgo} & \TableHeader{Ajuste aplicado} & \TableHeader{Efecto} \\ \hline
Sondeos ARP de DHCP & Se excluyó la IP origen \texttt{0.0.0.0} antes de construir conflictos ARP. & Eliminación de falsos conflictos asociados al arranque de teléfonos VoIP. \\ \hline
Alertas repetidas por HTTP & Se consolidaron flujos inseguros con la misma causa raíz en una alerta agregada. & Menor fragmentación del reporte y lectura más clara de severidad. \\ \hline
Falta de contexto de dispositivo & Se añadió el inventario \texttt{data/network/known\_hosts.json} para etiquetar MAC e IP reconocidas. & Alertas más legibles sin cambiar detecciones ni umbrales. \\ \hline
Impacto global & Reproceso de las 20 sesiones de ordeño posteriores a la intervención. & Alertas Alta por sesión: 31{,}6 antes del ajuste y 4{,}0 después ($-87\%$). \\ \hline
\end{tabularx}
\end{table}

La reducción de alertas Alta corresponde a ruido estructural del análisis, no a desaparición de todos los hallazgos. Las alertas que permanecen después del ajuste se relacionan con latencia elevada, eventos seriales incompletos, concurrencia en el bus y silencios prolongados del canal.

\section{Síntesis de cumplimiento de objetivos}

La evidencia presentada permite establecer el estado de cumplimiento de los cuatro objetivos específicos, con una distinción explícita entre resultados cerrados y resultados aún en consolidación estadística.

\begin{table}[H]
\centering
\caption{Síntesis de cumplimiento de objetivos.}
\label{tab:sintesis_cumplimiento_objetivos}
\footnotesize
\renewcommand{\arraystretch}{1.22}
\begin{tabularx}{\linewidth}{|>{\raggedright\arraybackslash}p{2.2cm}|>{\raggedright\arraybackslash}X|>{\centering\arraybackslash}p{2.4cm}|}
\hline
\rowcolor{gray!20} \TableHeaderFirst{Objetivo} & \TableHeader{Resultado principal} & \TableHeader{Estado} \\ \hline
Objetivo~1 & Se caracterizaron 63 visitas y 1\,818 sesiones; se identificaron conflictos ARP, exposición de telemetría en claro, fragmentación serial y baja visibilidad antes de la copia controlada de puerto. & Cumplido \\ \hline
Objetivo~2 & Se verificó una arquitectura perimetral con derivación serial aislada, captura pasiva, separación de dominios y evidencia preliminar de neutralidad eléctrica. & Cumplido \\ \hline
Objetivo~3 & La pasarela publicó sesiones con \ac{TLS}~1.3, sin retenciones ni fallos en el subconjunto consolidado, conservando el significado operativo generado por el motor. & Cumplido en el escenario evaluado \\ \hline
Objetivo~4 & La muestra previa a la intervención quedó cerrada ($n=64$, $\bar{\eta}=4.49\%$) y la posterior alcanzó el cierre planificado con $n=38$, $\bar{\eta}=22.80\%$, superando $n\geq30$. El contraste Mann-Whitney~U confirma diferencia significativa. & Cumplido \\ \hline
\end{tabularx}
\end{table}

En conjunto, los resultados respaldan la pertinencia de la intervención perimetral: la solución mejora la visibilidad, ordena la publicación de eventos y reduce la pérdida silenciosa de evidencia normalizada. La conclusión estadística del Objetivo~4 queda cerrada con el contraste Mann-Whitney~U ejecutado sobre $n_{\text{pre}}=64$ y $n_{\text{post}}=38$, confirmando una diferencia significativa de eficiencia de extracción a favor del escenario instrumentado.
